\chapter{星海的记忆载体}

巨鲸内腔消失后，投送港的金属地面还在洛安脚下轻轻震动。

那震动并不来自平台。它像从掌心那枚种子里传出来，顺着骨头往上爬，一路爬进胸腔。洛安低头看着自己的手，手指仍在发抖。查理也在发抖，只不过他没有真正的手，只能跟着这具身体一起承受。

“你看见了多少？”瑟恩问。

他的声音比刚才更哑，翼膜边缘垂在身后，像两片被火烧过的透明布。

洛安张了张嘴，却一时说不出话。

他看见的东西太多了。恒星的诞生、文明的迁徙、无数不属于任何已知种族的面孔，还有一些完全无法用“面孔”来称呼的生命形式。它们有的生活在气体海里，有的把身体展开成环绕母星的膜，有的只剩下光和节律。巨鲸从它们身边经过，收走那些即将断掉的声音。

不是抢夺。

查理立刻在心里停了一下，把这个句式吞回去。

巨鲸经过那些文明时，更像一艘沉默的船靠岸。那些即将毁灭的族群把自己的最后一部分交给它：一段航线，一种语言，一个失败的原因，或者一个孩子死前唱过的旋律。

洛安终于说：“它们装着很多死去的世界。”

渊民族老人抬起头，神情肃穆：“每一头巨鲸都承载一段无法再靠自己延续的历史。有的背负一个族群，有的背负一条航线，有的只背负一个错误。”

“错误也要保存？”索勒问。

老人看向他：“尤其是错误。胜利者总喜欢把错误改成荣耀，失败者又常常来不及把原因说完。巨鲸保存它们，是为了让后来者知道，文明并非只因为敌人而死。”

岚柯把罐体抱得更紧。她没有反驳，眼神却落在那些正在崩裂的轨道残骸上。

投送港外，法厄同的天空已经失去稳定颜色。紫色与黑色混在一起，远方有一整段地壳缓缓抬起，像被星球内部的手翻开。碎片在低空形成一圈缓慢旋转的带状阴影。

查理知道，那些碎片会在很久以后变成小行星带的一部分。可洛安不知道。他只觉得天空在往下压，时间也在压。

巨鲸种子又亮了一次。

这一次，洛安没有进入完整幻象，只看见几段短促的画面。

第一段，是一头巨鲸穿过木星外缘。它的身体表面有无数暗蓝色纹路，每一道纹路都对应着一段记忆。它没有舰队护送，也没有王座仪仗，只有十几艘小型飞行器远远跟着，像在护送一座移动的墓园。

第二段，是某个无名文明的最后一天。天空里有三颗太阳，其中两颗已经熄灭。数百个透明生命围在一片冰原上，把自己的语言压缩成光点，送进巨鲸打开的腔室。

第三段，是法厄同战争前夜。五族代表在中立环室里争执，巨鲸停在环室外的黑暗中，庞大得像一段夜色。它没有发言，只把每一句威胁、每一次沉默、每一个退让的失败都记录下来。

洛安忽然明白，巨鲸最强大的能力并非毁灭。

它让历史无法只剩下胜利者的版本。

这个念头像一块冷石头落进胸口。查理也被砸了一下。他想到自己时代那些被写进教材、纪念馆和游戏里的战争。每一场战争都有一个更容易讲述的版本，有英雄，有敌人，有结束。可真正活过那一天的人，可能只记得烟、血、跑不动的腿，还有自己为什么没能拉住旁边的人。

“所以它们一直在看着我们？”洛安问。

“看着，听着，记着。”渊民族老人说，“但它们很少替谁做决定。”

“这次做了。”

老人沉默。

远处传来爆炸，平台一侧的环形支架折断，数十具尸体和残骸被抛向外壳裂缝。索勒伸手抓住一根断柱，岚柯也几乎摔倒。瑟恩靠最后一点翅振稳住身体，脸色白得可怕。

“路快没了。”瑟恩说。

洛安把巨鲸种子按回胸前的临时封袋里。

那一瞬间，他感觉自己像把一整个星海塞进一块破布里。

这太荒唐了。

也太重了。

可脚下的光路还在，尽头通向投送港深处最后一艘残舟。

洛安咬住牙，迈出一步。
